\documentclass[12pt]{article}
\usepackage[utf8]{inputenc}
\usepackage[T1]{fontenc}
\usepackage[margin=1in]{geometry}
\usepackage{setspace}
\usepackage{xcolor}
\usepackage[colorlinks=true, urlcolor=blue!60!black, linkcolor=blue!60!black]{hyperref}
\setlength{\parindent}{0em}
\setlength{\parskip}{0.45em}
\renewcommand{\baselinestretch}{1.0}
\begin{document}
\begin{center}
{\Large \textit{Prigg v. Pennsylvania}, 41 U.S. 539 (1842)\par}
\vspace{0.3em}
{\normalsize Case Brief --- One-Page Edition\par}
\end{center}

\textbf{Memory Jogger.} Under the Fugitive Slave Clause, slaveholders had a self-executing right to recapture fugitives, and Congress's rendition power was exclusive (\href{https://supreme.justia.com/cases/federal/us/41/539/}{Justia}).

\textbf{Facts.} Edward Prigg forcibly carried Margaret Morgan and her children from Pennsylvania to Maryland, where her owner held her as a slave, violating Pennsylvania's 1826 anti-kidnapping statute.

\textbf{Procedural History.} In an agreed test case, Prigg was convicted on a special verdict; affirmed \textit{pro forma}.

\textbf{Votes.} Reversed 8--1: Story for the Court; four concurrences in the judgment; McLean dissented.

\textbf{Holding.} State laws obstructing peaceable recaption are void; federal power is exclusive; States need not enforce the federal Act.

\textbf{Reasoning.} The Clause created a ``positive, unqualified right,'' and Congress's implied enforcement power bars even ``auxiliary'' state laws. State magistrates may act only ``unless prohibited by state legislation.'' Taney would have let States aid slaveholders; McLean would have required a federal certificate before removal.

\textbf{Examples.} \emph{Same-Side:} extra state rules on federally regulated railroads. \emph{Opposite-Side:} withdrawing state officers' help. \emph{Fence-Sitter:} state penalties ``in aid'' of federal rules.

\textbf{Critique.} Finkelman shows Story stripped free Black Northerners of protection against kidnapping.

\textbf{Key Quotation.} ``Upon this ground, we have not the slightest hesitation in holding that, under and in virtue of the Constitution, the owner of a slave is clothed with entire authority, in every State in the Union, to seize and recapture his slave whenever he can do it without any breach of the peace or any illegal violence. In this sense and to this extent, this clause of the Constitution may properly be said to execute itself, and to require no aid from legislation, state or national.'' --- Justice Story, 41 U.S. at 613.

\end{document}
